\documentclass{article}

\usepackage[utf8]{inputenc}
\usepackage[spanish]{babel}
\usepackage{amssymb}
\usepackage{graphicx}
\usepackage{amsbsy}
\usepackage{pifont}
\usepackage{epsfig}

\begin{document}

\hyphenation{}

\title{Aplicación de Inteligencia Artificial para la Interacción y Análisis de Datos del Proyecto “Semáforo de Pobreza” en Paraguay}

\author{
  Carlos Bustamante \\
  Universidad Comunera \\
  \texttt{cbustamantem@gmail.com}
  \and
  Nelba Barreto \\
  Universidad Comunera \\
  \texttt{barretonelba@gmail.com}
}

\date{Noviembre 2025}
\maketitle

\section{Introducción}\label{sec:introduccion}
La pobreza es un fenómeno complejo que afecta a millones de personas alrededor del mundo, y Paraguay no es la excepción. En este contexto, el proyecto Semáforo de Pobreza (Poverty Stoplight), desarrollado por la Fundación Paraguaya, se presenta como una herramienta innovadora que permite a las familias autoevaluar su situación en diversas dimensiones y construir un plan de acción personalizado para superar sus limitaciones.

Actualmente, el Semáforo de Pobreza cuenta con más de 708.000 evaluaciones realizadas en 59 países y una base de conocimientos de más de 1.000 soluciones prácticas implementadas por distintas comunidades. Sin embargo, el acceso y aprovechamiento de esta vasta información sigue siendo limitado, tanto para nuevas organizaciones interesadas como para las personas interesadas en participar en el programa.

En este sentido, surge la necesidad de construir un agente social basado en inteligencia artificial, capaz de interactuar con personas y organizaciones de manera natural y ofrecer respuestas informadas, ejemplos de casos de éxito, estadísticas y estrategias efectivas extraídas de la base de conocimientos del Semáforo. La aplicación de técnicas modernas de Retrieval Augmented Generation (RAG) y fine tunning de modelos de lenguaje (LLM) permitirá estructurar, comprender y comunicar esta información de forma accesible, confiable y contextualizada.

\section{Estado del arte}\label{sec:estadoArte}
A nivel global, diferentes proyectos han buscado aplicar inteligencia artificial en el ámbito del desarrollo social y la reducción de la pobreza. Iniciativas como AI for Good (ONU) y plataformas de análisis de datos sociales han demostrado el potencial de los modelos de lenguaje para sintetizar conocimiento, generar reportes y apoyar intervenciones sociales.

Sin embargo, la mayoría de los esfuerzos se concentran en países desarrollados o se basan en datos estructurados limitados. En el caso del Semáforo de Pobreza, la información disponible es vasta, pero heterogénea y distribuida entre múltiples bases de datos, lo que dificulta su uso para generar conocimiento aplicable.

Hasta la fecha, no existe un modelo de inteligencia artificial entrenado específicamente con los datos y soluciones del Semáforo de Pobreza, ni un agente conversacional que facilite la transferencia de estas experiencias a nuevas comunidades y organizaciones. Por tanto, esta investigación busca cubrir esa brecha mediante un enfoque integral que combine ciencia de datos, procesamiento de lenguaje natural y aprendizaje automático.

\section{Descripción de la propuesta}\label{sec:descripcion}
La propuesta consiste en diseñar y desarrollar un agente social inteligente basado en técnicas de RAG (Retrieval-Augmented Generation), que integre toda la base de conocimientos del Semáforo de Pobreza y permita interactuar de manera conversacional con usuarios y organizaciones interesadas en aplicar el semáforo de pobreza.

La idea principal es aplicar modelos de lenguaje (LLM) entrenados con datos reales del Semáforo —incluyendo indicadores, planes de acción, casos de éxito, estadísticas por país, y la base de conocimientos la cual consiste en más de 1000 casos de exitos, información actualmente registrada en su plataforma interna. — para ofrecer respuestas útiles y contextualizadas. Este agente podrá responder preguntas como:
“¿Cuáles fueron los resultados del Semáforo en una cadena de supermercados en Paraguay?”
y mostrar, por ejemplo, que durante el proceso se detectaron un alto porcentaje de sus funcionarios tenían carencias en vivienda, acceso a energía y transporte, junto con las soluciones implementadas y los resultados luego de la asistencia.

Este sistema contribuirá a facilitar el conocimiento social acumulado por el Semáforo y fortalecer la toma de la interacción con organizaciones y personas interesadas.

\section{Objetivos}\label{sec:objetivos}
\begin{itemize}
 \item \textbf{Objetivo general}
   \begin{itemize}
     \item Desarrollar un agente social basado en inteligencia artificial que integre la base de conocimientos del proyecto Semáforo de Pobreza, utilizando técnicas de ciencia de datos y modelos de lenguaje para facilitar la interacción y difusión de experiencias exitosas.
   \end{itemize}
 \item \textbf{Objetivos Específicos}
   \begin{itemize}
    \item Desarrollar y estructurar la base de conocimientos del Semáforo de Pobreza consolidando sus diversas fuentes de datos.
    \item Construir un modelo de lenguaje (LLM) entrenado con información del Semáforo y técnicas RAG.
    \item Implementar una interfaz web interactiva que permita la comunicación con el agente social inteligente en español e inglés.
    \item Demostrar la precisión y relevancia de las respuestas generadas por el modelo.
    \item Evaluar la utilidad del sistema en distintos escenarios de uso por parte de organizaciones y usuarios.
   \end{itemize}
\end{itemize}

\section{Metas}\label{sec:metas}
Las metas de la propuesta se pueden enumerar de la siguiente manera:
\begin{enumerate}
    \item Integrar en un único repositorio estructurado la base de conocimientos actual del Semáforo.
    \item Entrenar un modelo de lenguaje de dominio específico.
    \item Desarrollar un prototipo funcional del agente social.
    \item Validar el sistema mediante pruebas con usuarios reales y organizaciones asociadas.
    \item Publicar los resultados obtenidos y su impacto social.
\end{enumerate}

\section{Aportación científica esperada}\label{sec:aporte}
Al finalizar el proyecto se contará con:
\begin{enumerate}
\item Un repositorio en Hugging Face con la base de conocimientos utilizada para realizar el fine tuning del modelo de lenguaje (LLM) especializado en datos de desarrollo social y pobreza.
\item Un método estructurado para integrar y aprovechar grandes bases de conocimiento sociales.
\item Un agente conversacional funcional, accesible desde el portal del Semáforo de Pobreza povertystoplight.org.
\item Un avance científico y tecnológico en el uso de IA aplicada a políticas públicas y proyectos de impacto social.
\end{enumerate}

\section{Metodología científica}\label{metodologia}
En la etapa inicial, se realizará un trabajo de ciencia de datos enfocado en recopilar, limpiar y estructurar la base de conocimientos del Semáforo de Pobreza.

Seguidamente, se entrenará un modelo LLM con técnicas de RAG, utilizando los datos estructurados y textos descriptivos de las soluciones históricas.

Finalmente, se desarrollará una interfaz web que permitirá la interacción con el agente, validando su rendimiento con métricas de precisión, relevancia y satisfacción del usuario.

\section{Calendario de actividades}\label{sec:calendario}
El calendario de actividades se divide en las siguientes fases es y están detalladas en la Figura xx:\\

Fase 1: Preparación y estructuración de la base de conocimiento (1 mes)
Fase 2: Entrenamiento del modelo LLM con la base de conocimiento (1 mes)\\
Fase 3: Implementación del agente virtual en el sitio web Poverty Stoplight (1 mes)\\
Fase 4: Pruebas piloto, evaluación de resultados y ajustes finales (1 mes)\\

Aquí introducir diagrama de gantt con el calendario por fases o por actividades.
%\begin{figure}
%\includegraphics[width=13cm, height=7cm]{diagramaGant}
%\caption{Calendario de Actividades}\label{fig:calendario}
%\end{figure}


\bibliographystyle{plain}
\bibliography{protocolo}

\clearpage
% Visto bueno
\vfill
\Large{Vo. Bo.}
\\    \\    \\    \\    \\  \\    \\    \\    \\    \\   
 
\begin{tabular}{@{}p{6cm}p{7cm}@{}}
    \Large{Firma del Alumno:} & \hrulefill \\
                         & \large{Alumno} \par \\
    \\    \\    \\    \\    \\    
    \Large{Directora de tesis:} & \hrulefill \\
                         & \large{Dra. Helena Gómez Adorno} \\
    \Large{} & \\
                         & \large{Investigadora Titular A} \\
    \Large{} &  \\
                         & \large{IIMAS-UNAM} \\
    
\end{tabular}


\end{document}
